% GENERATED FILE. Edit canonical lessons, then run npm run generate:books.
% canonical-source-hash: fnv1a64:8d245652b097aa3c
% canonical-lessons: FR-C259-dabord, FR-C259-raconter, FR-C259-lecture, FR-C259-revision-histoire

\chapter{D'abord, Ensuite, Enfin — A Story Told in Order}
\label{ch:fr-a-story-told-in-order}
\hlchaptermodality{259}

\begin{chapteropening}
\textbf{By the end of this chapter:} \emph{I can tell a short connected story in order with d'abord, ensuite and enfin, and read one of about a hundred words.}

Complete FR-C259-revision-histoire: retell the reading's journey aloud, choosing scene or event for every verb.
\end{chapteropening}

\section[d'abord]{d'abord — the events, in order}

\label{lesson:FR-C259-dabord}



\begin{quote}
Say \textquotedblleft{}then\textquotedblright{}. (\emph{Puis.}) Then \textquotedblleft{}afterwards\textquotedblright{}. (\emph{Ensuite.}) Then \textquotedblleft{}finally\textquotedblright{}. (\emph{Enfin.})

Three words for the middle and the end of a story. The order still needs a first word.
\end{quote}

\subsection*{You'll want to know: d'abord}
\textbf{d'abord} — \textquotedblleft{}first\textquotedblright{}, \textquotedblleft{}at first\textquotedblright{}.

\begin{quote}
\textbf{D'abord}, j'ai acheté le billet. — First I bought the ticket. \textbf{Ensuite}, j'ai attendu le train. — Next I waited for the train. \textbf{Puis} le train est arrivé. — Then the train came. \textbf{Enfin}, je suis rentré à la maison. — Finally I went back home.
\end{quote}

Say it as one word: \textquotedblleft{}da-BOR\textquotedblright{}. The \emph{d} is \emph{de} with its vowel dropped, and the final \emph{d} of \emph{abord} is silent.

\begin{grammarlens}[title={Grammar Lens: four places in a story}]
\textbf{Ensuite} and \textbf{puis} both mean \textquotedblleft{}then, next\textquotedblright{}, and a story can use them as many times as it has events. \emph{Puis} is the lighter of the two, and speech often doubles it up as \emph{et puis}.

\textbf{Enfin} closes the story, and it often carries relief: something you were waiting for has come at last. In speech it has a second life as a hesitation, \emph{enfin…}, \textquotedblleft{}well, I mean…\textquotedblright{}, which you will hear constantly.

\textbf{D'abord} can also mean \textquotedblleft{}at first\textquotedblright{}, setting up a change:

\begin{quote}
\emph{D'abord, j'avais peur, mais ensuite, c'était bien.} — At first I was scared, but then it was fine.
\end{quote}

Notice what these words do \textbf{not} touch: the scene. \emph{Il faisait froid} needs no \emph{d'abord}. Order belongs to the events.
\end{grammarlens}

\begin{cousinweb}
\textbf{d'abord} is \emph{de} plus \textbf{abord}, \textquotedblleft{}the approach, the first coming alongside\textquotedblright{}. \emph{Abord} is built on \textbf{bord}, \textquotedblleft{}the edge, the side of a ship\textquotedblright{}, a Germanic word: EtymDB 2.1 sets it beside Old English \emph{bord}, \textquotedblleft{}side of a ship\textquotedblright{}, and Old Norse \emph{borð}. The usual account runs through the verb \emph{aborder}, to bring a ship alongside: \emph{d'abord} is \textquotedblleft{}from the very first approach\textquotedblright{}.
\end{cousinweb}

\subsection*{Your turn}
Say these aloud:

\begin{itemize}
\raggedright
  \item \emph{D'abord, j'ai acheté le billet. Ensuite, j'ai attendu le train.}
  \item \emph{Puis le train est arrivé. Enfin, je suis rentré à la maison.}
  \item at first, then a change — \emph{D'abord, j'avais peur, mais ensuite, c'était bien.}
  \item a scene, then events — \emph{Il pleuvait. D'abord, j'ai attendu. Ensuite, quand le train est arrivé, je suis monté.}
\end{itemize}

\begin{tcolorbox}[breakable,colback=teal!4,colframe=teal!35!black,title={Before you move on}]
Which word opens a story told in order? (\emph{\textbf{D'abord.}}) Which two words mean \textquotedblleft{}then\textquotedblright{}? (\emph{\textbf{Ensuite}} and \emph{\textbf{puis}}.) Does the weather need a \emph{d'abord}? (\textbf{No} — order belongs to the events.) And what was \emph{abord}, at root? (A ship's \textbf{first coming alongside}.)
\end{tcolorbox}
\section[raconter]{raconter — to tell a story, which is to count it}

\label{lesson:FR-C259-raconter}



\begin{quote}
You can hold a scene still and put events in order with \emph{d'abord}. This lesson gives the act a name, then puts it all together.
\end{quote}

\subsection*{You'll want to know: raconter}
\textbf{raconter} — \textquotedblleft{}to tell\textquotedblright{}, for a story or for what happened. It is a regular \emph{-er} verb: \emph{je raconte}, \emph{j'ai raconté}, \emph{je racontais}.

\begin{quote}
\emph{Je raconte une histoire à mon frère.} — I tell my brother a story. \emph{Raconte !} — Go on, tell me!
\end{quote}

You have \textbf{l'histoire}, which is both \textquotedblleft{}history\textquotedblright{} and \textquotedblleft{}a story\textquotedblright{}, so \textbf{raconter une histoire} is the whole act in three words. The person who hears it comes after \textbf{à}.

\begin{grammarlens}[title={Grammar Lens: the backbone of a told story}]
Here is what a story tells, with both pasts to work with:

\begin{quote}
\textbf{C'était} samedi. \textbf{Il faisait} beau. \textbf{D'abord, je suis allé} au marché. \textbf{Ensuite, j'ai rencontré} un ami, et nous \textbf{avons parlé}. \textbf{Mais il a commencé} à pleuvoir, et \textbf{enfin je suis rentré} à la maison.
\end{quote}

\begin{quote}
\emph{It was Saturday. The weather was fine. First I went to the market. Then I met a friend, and we talked. But it started to rain, and in the end I went back home.}
\end{quote}

\emph{C'était}, \emph{il faisait}: the scene, held still. \emph{Je suis allé}, \emph{j'ai rencontré}, \emph{nous avons parlé}, \emph{il a commencé}, \emph{je suis rentré}: the events, put in order by \emph{d'abord}, \emph{ensuite} and \emph{enfin}. \textbf{Mais} marks the turn, the point where the story stops going the way it was going. \emph{Et} adds; \emph{mais} turns.
\end{grammarlens}

\begin{cousinweb}
\textbf{raconter} goes back to Old French \emph{raconter} (EtymDB 2.1). It is built on \textbf{conter}, \textquotedblleft{}to tell\textquotedblright{}, which EtymDB 2.1 traces to Latin \emph{computāre}, \textquotedblleft{}to reckon up, to compute\textquotedblright{} (Lewis \& Short, under \emph{computo}), with \emph{re-} and \emph{a-} in front. To tell a story was to count it out, one piece after another. English kept both senses in one word: you \textbf{count} a sum, and you \textbf{recount} what happened.
\end{cousinweb}

\subsection*{Your turn}
Say these aloud:

\begin{itemize}
\raggedright
  \item the scene — \emph{C'était samedi. Il faisait beau.}
  \item the events — \emph{D'abord, je suis allé au marché. Ensuite, j'ai rencontré un ami.}
  \item the turn and the end — \emph{Mais il a commencé à pleuvoir, et enfin je suis rentré à la maison.}
  \item \emph{Je raconte une histoire à mon frère.}
\end{itemize}

\begin{tcolorbox}[breakable,colback=teal!4,colframe=teal!35!black,title={Before you move on}]
In your story, which past held the weather and which one took you to the market? (\emph{\textbf{Il faisait}} — imparfait; \emph{\textbf{je suis allé}} — passé composé.) Which word marked the turn? (\emph{\textbf{Mais.}}) Which Latin verb is behind \emph{conter}? (\emph{\textbf{Computāre}}, \textquotedblleft{}to reckon up\textquotedblright{}.)
\end{tcolorbox}
\section[(reading)]{(reading) — a journey, told in both pasts}

\label{lesson:FR-C259-lecture}



\begin{quote}
Chapter forty-five gave you sixty-three words in a row, all in the present. This passage is longer, and it happened last Saturday.

Every word in it has come up before this lesson.
\end{quote}

\begin{tcolorbox}[breakable,skin=enhanced,colback=orange!6,colframe=orange!45!black,boxrule=0.5pt,arc=1mm,left=6pt,right=6pt,top=4pt,bottom=4pt,fonttitle=\bfseries,title={Read this}]
\begin{quote}
Samedi dernier, c'était l'anniversaire de mon ami. Il faisait froid et il pleuvait. D'abord, j'ai attendu le train sur le quai. Il y avait beaucoup de monde, mais le train n'arrivait pas, et j'étais inquiet. Ensuite, le train est arrivé, et je suis monté.
\end{quote}

\begin{quote}
Dans le train, je lisais, et un enfant dormait à côté de moi. Quand je suis arrivé, mon ami m'attendait devant la maison. Puis nous avons mangé du gâteau, et nous avons beaucoup parlé. Enfin, le soir, je suis rentré à la maison. Il faisait beau, et il y avait des étoiles. J'étais content.
\end{quote}

Read it once without stopping. Do not translate as you go.

Then read it a second time and sort every verb into one of two piles: the scene, or what happened.
\end{tcolorbox}

\begin{culture}
The events pile is short: \emph{j'ai attendu}, \emph{est arrivé}, \emph{je suis monté}, \emph{je suis arrivé}, \emph{nous avons mangé}, \emph{nous avons parlé}, \emph{je suis rentré}. Seven things happened, and \emph{d'abord}, \emph{ensuite}, \emph{puis} and \emph{enfin} put them in order.

Everything else is scene: the birthday, the cold, the rain, the crowd, the worry, the stars. A story is mostly background.

\emph{Le train n'arrivait pas} is scene: the not-arriving went on. \emph{Le train est arrivé} is the event that ended it. And \emph{mon ami m'attendait} is scene too. He was already waiting, and the narrator's arrival cut into it, with \emph{quand}.
\end{culture}

\subsection*{Your turn}
\begin{itemize}
\raggedright
  \item \emph{Say it:} the passage aloud, once, without stopping
\end{itemize}

\begin{tcolorbox}[breakable,colback=teal!4,colframe=teal!35!black,title={Before you move on}]
How many words in that passage were new? (\textbf{None.})

Why is \emph{m'attendait} in the imparfait? (The waiting \textbf{was already going on} when the narrator arrived: it is scene.)

Which word opened the line where a scene met an event? (\emph{\textbf{Quand.}})
\end{tcolorbox}
\section[(review)]{Review — the journey, retold}

\label{lesson:FR-C259-revision-histoire}



\begin{quote}
Say \textquotedblleft{}first\textquotedblright{}. (\emph{D'abord.}) Say \textquotedblleft{}I tell a story\textquotedblright{}. (\emph{Je raconte une histoire.})
\end{quote}

\begin{grammarlens}[title={What you've built}]
Two chapters, and a whole story. A scene in the imparfait: \emph{il faisait froid}, \emph{il y avait beaucoup de monde}, \emph{j'étais inquiet}. Events in the passé composé, held in line by \emph{d'abord}, \emph{ensuite}, \emph{puis} and \emph{enfin}. A turn with \emph{mais}. And the joint where a scene meets an event: \emph{quand}.

Now the journey from the reading, told by you. Here is its skeleton, scene first, then the events:

\begin{quote}
\emph{C'était samedi. Il faisait froid et il pleuvait.} \emph{D'abord, j'ai attendu le train. Ensuite, je suis monté.} \emph{Quand je suis arrivé, mon ami m'attendait.} \emph{Enfin, je suis rentré à la maison. J'étais content.}
\end{quote}
\end{grammarlens}

\subsection*{Your turn}
Say these aloud:

\begin{itemize}
\raggedright
  \item the scene — \emph{C'était samedi. Il faisait froid et il pleuvait. Il y avait beaucoup de monde.}
  \item the first two events — \emph{D'abord, j'ai attendu le train. Ensuite, je suis monté.}
  \item a scene inside the journey — \emph{Pendant que je lisais, un enfant dormait.}
  \item the event that met the waiting friend — \emph{Quand je suis arrivé, mon ami m'attendait.}
  \item the end — \emph{Enfin, je suis rentré à la maison. J'étais content.}
\end{itemize}

\begin{tcolorbox}[breakable,colback=teal!4,colframe=teal!35!black,title={Before you move on}]
Which word puts the first event in place? (\emph{\textbf{D'abord.}}) Is \emph{j'étais content} scene or event? (\textbf{Scene} — a feeling has no edge.) Which word marks the turn in a story? (\emph{\textbf{Mais.}}) And which Latin verb is inside \emph{raconter}? (\emph{\textbf{Computāre}}, to count it out.)
\end{tcolorbox}
